% Generated draft: requires mathematical and source review.
\documentclass[11pt]{article}
\usepackage[margin=25mm]{geometry}
\usepackage{fontspec}
\setmainfont{Latin Modern Roman}
\usepackage{amsmath,amssymb,mathtools,unicode-math}
\setmathfont{Latin Modern Math}
\usepackage{enumitem,needspace,xurl,hyperref}
\setlength{\parindent}{0pt}
\setlength{\parskip}{5pt}
\setlength{\emergencystretch}{3em}
\newcommand{\N}{\mathbb{N}}
\newcommand{\Z}{\mathbb{Z}}
\newcommand{\Q}{\mathbb{Q}}
\newcommand{\R}{\mathbb{R}}
\newcommand{\C}{\mathbb{C}}
\newcommand{\F}{\mathbb{F}}
\newcommand{\A}{\mathbb{A}}
\newcommand{\PP}{\mathbb{P}}
\newcommand{\E}{\mathbb{E}}
\newcommand{\eps}{\varepsilon}
\newcommand{\dd}{\,\mathrm{d}}
\begin{document}
\section*{UnsolvedMath ID: 8700016}
\textbf{Source problem code:} catalogue entry
\subsection{Definitions and mathematical statement}
Let \(\theta\in\R\) be algebraic over \(\Q\) with \([\Q(\theta):\Q]\geq3\). Waldschmidt's Conjecture 2.12 proposes the equivalence
\[
 \begin{gathered}
 \text{inequality (2.11) has infinitely many solutions }(p,q)\in\Z\times\Z_{>0}
 \end{gathered}
 \quad\Longleftrightarrow\quad
 \int_1^\infty\psi(x)\,\dd x\text{ diverges}.
\]
Here inequality (2.11) and the assumptions on \(\psi\) are not supplied. Consequently its exact error bound, strictness, normalization in \(q\), admissible regularity or monotonicity of \(\psi\), and any coprimality condition are unknown. The integral is an improper integral under whatever hypotheses the missing source context imposes. In particular its divergence cannot be identified with \(+\infty\) without the intended nonnegativity assumption.

A standard-looking approximation bound such as \(|\theta-p/q|<\psi(q)/q\) must not be substituted for the absent (2.11) merely because the summary mentions Khinchin-type questions: using \(\psi(q)\), \(\psi(q)/q\), or \(\psi(q)/q^2\) would change the statement. The conjecture is pointwise for each specified algebraic \(\theta\), not a Lebesgue-almost-everywhere claim over real numbers. Both directions of the integral criterion are retained, pending recovery of the missing inequality and function class.
\subsection{Short English statement}
For a fixed real algebraic number of degree at least three, the conjecture gives an integral divergence criterion for infinitely many rational approximations. The defining approximation inequality and the hypotheses on its function \(\psi\) are missing.
\subsection{Sources}
\begin{itemize}
\item Local UnsolvedMath registry, permanent ID 8700016, source code catalogue entry; attributed to Michel Waldschmidt; no proposal year supplied; updated 2026-08-06.
\item Waldschmidt, \emph{Open Diophantine Problems} (2004), Conjecture 2.12: \url{https://arxiv.org/abs/math/0312440}. Registry access date 2026-07-29; source-TeX extraction. The source's numbered inequality (2.11) is referenced but not included. Current primary-literature status and redistribution terms require review.
\item The supplied summary compares the question to Khinchin and Duffin--Schaeffer theory and Roth's theorem, without a precise additional citation or a verified resolution of this fixed-algebraic-number claim.
\end{itemize}
\end{document}
